\chapter{Internationalization System}
\label{ch:i18n-system}

OmniLaTeX provides a layered internationalization (i18n) system that handles
language rules, string localization, and bidirectional text.  This chapter
covers the translation infrastructure; see Chapters~41 and~42 for multilingual
documents and RTL scripts respectively.

% =========================================================================
\section{polyglossia Setup}
% =========================================================================

OmniLaTeX loads \texttt{polyglossia} through \texttt{omnilatex-i18n.sty},
which serves as the LuaLaTeX replacement for \texttt{babel}.  The primary
document language is set via the class option \texttt{language=<lang>} and
propagated to \texttt{polyglossia} using a double-\cs{expandafter} pattern
that ensures the macro is fully expanded before \cs{setdefaultlanguage}
receives it.

In addition to the primary language, 27 secondary languages are registered
via \cs{setotherlanguages}, making them available for inline use with
\cs{begin}\texttt{\{<lang>\}}\ldots\cs{end}\texttt{\{<lang>\}} or
\cs{foreignlanguage}\texttt{\{<lang>\}\{...\}}:

\begin{verbatim}
\setotherlanguages{
  german, ngerman, english, french, spanish, italian, portuguese,
  russian, dutch, polish, czech, greek, turkish, simplifiedchinese,
  traditionalchinese, japanese, korean, arabic, hebrew, vietnamese,
  hindi, swedish, finnish, danish, norsk
}
\end{verbatim}

This design allows a German thesis to include English passages with a simple
\cs{begin}\texttt{\{english\}}\ldots\cs{end}\texttt{\{english\}} switch,
or an English document to embed Arabic text inline.

% =========================================================================
\section{Translation System}
% =========================================================================

Beyond hyphenation and caption rules, OmniLaTeX uses the \texttt{translations}
package for localized strings that \texttt{polyglossia} does not cover
(e.g.\ custom headings, institutional labels).  Translations are declared
with \cs{DeclareTranslation} and retrieved with \cs{GetTranslation}:

\begin{verbatim}
\DeclareTranslation{german}{Gloss}{Glossar}
\GetTranslation{Gloss}   % → "Glossary" or "Glossar" depending on locale
\end{verbatim}

Users can override any translation without modifying the package:

\begin{verbatim}
\RenewTranslation{german}{author}{Autorin}
\end{verbatim}

This is particularly useful for gender-neutral or gender-specific academic
roles.  The document preamble in \texttt{config/document-settings.sty}
demonstrates overriding \texttt{author}, \texttt{FirstExaminer},
\texttt{SecondExaminer}, and \texttt{Supervisor} for German documents.
OmniLaTeX also uses \cs{GetTranslationFor}\texttt{\{<lang>\}\{<key>\}} to
inspect translations for a specific language, enabling conditional logic such
as selecting the correct German definite article (\texttt{der}/\texttt{die})
based on the gendered form of \texttt{author}.

% =========================================================================
\section{Language Fallback Chain}
% =========================================================================

When a translation key is not defined for the current document language,
the \texttt{translations} package falls back to English.  This behaviour is
formally verified in \texttt{LanguageFallback.lean}, which defines the
fallback function as:

\begin{verbatim}
def fallback : PrimaryLanguage → PrimaryLanguage
  | .english => .english
  | _        => .english
\end{verbatim}

Two properties are machine-checked:
\begin{itemize}
  \item \texttt{fallback\_terminates}: for every language~$l$,
    $\text{fallback}(l) = \text{english}$.
  \item \texttt{fallback\_stable}:
    $\text{fallback}(\text{fallback}(l)) = \text{fallback}(l)$,
    guaranteeing that repeated fallbacks do not oscillate.
\end{itemize}

English also serves as the fixed point: $\text{fallback}(\text{english})
= \text{english}$.  This ensures that all 18 primary languages resolve
to a defined string even if a particular key is missing.

% =========================================================================
\section{Dynamic Language Switching}
% =========================================================================

Any registered secondary language can be activated inline using
polyglossia's environment form:

\begin{verbatim}
\begin{english}
  This paragraph uses English hyphenation and caption rules.
\end{english}
\end{verbatim}

This is essential for bilingual documents where, for example, a German thesis
must include an English abstract for international readers, or where legal
disclaimers must appear in a specific language.  The language switch affects
hyphenation patterns, caption strings, and \cs{GetTranslation} resolution
for the duration of the environment.

% =========================================================================
\section{Translation Key Categories}
% =========================================================================

The 47 translation keys are organized into functional categories.
Completeness is verified in \texttt{I18nCompleteness.lean}:

\begin{description}
  \item[Document metadata (7)] \texttt{Title}, \texttt{author},
    \texttt{CompiledOn}, \texttt{CensorNotice}, \texttt{PlaceDate},
    \texttt{Topic}, \texttt{Task}, \texttt{For}.
  \item[Academic roles (7)] \texttt{FirstExaminer},
    \texttt{SecondExaminer}, \texttt{Supervisor}, \texttt{Examiner},
    \texttt{AuthorshipDeclTitle}, \texttt{AuthorshipDeclText},
    \texttt{authorArticle}.
  \item[Glossary groups (10)] \texttt{Gloss}, \texttt{Greek},
    \texttt{Roman}, \texttt{Other}, \texttt{Terms}, \texttt{Names},
    \texttt{Acronyms}, \texttt{Subscripts}, \texttt{Superscripts},
    \texttt{SubSuperTitle}.
  \item[Content lists (10)] \texttt{ListOfExamples},
    \texttt{ListOfListings}, \texttt{ListOfIllustrations},
    \texttt{IllustrativeElements}, \texttt{Listing}, \texttt{Listings},
    \texttt{Example}, \texttt{Examples}, \texttt{FurtherReadingTitle},
    \texttt{FurtherReadingText}.
  \item[Scientific notation (4)] \texttt{Constant}, \texttt{Unit},
    \texttt{PhysicalQuantity}, \texttt{AdaptedFrom}.
  \item[Reactions (2)] \texttt{Reaction}, \texttt{Reactions}.
  \item[Table formatting (2)] \texttt{Contfoot}, \texttt{Conthead}.
  \item[Layout (5)] \texttt{BlankPage}, \texttt{First}, \texttt{Second},
    \texttt{LatexClass}, \texttt{Generator}.
\end{description}

% =========================================================================
\section{Complete Key Table}
% =========================================================================

The following table lists representative keys from each category.
The full list of all 47 keys with translations in all 18 languages is
documented in Chapter~44.

\begin{table}[htbp]
  \centering
  \caption{Selected translation keys and their English values.}
  \label{tab:translation-keys}
  \begin{tabular}{@{}lll@{}}
    \toprule
    Category & Key & English value \\
    \midrule
    Metadata & \texttt{CompiledOn}     & Compiled on \\
    Metadata & \texttt{PlaceDate}       & Place \& Date \\
    Academic & \texttt{FirstExaminer}   & Examiner \\
    Academic & \texttt{Supervisor}      & Supervisor \\
    Academic & \texttt{AuthorshipDeclTitle} & Declaration of Authorship \\
    Glossary & \texttt{Gloss}           & Glossary \\
    Glossary & \texttt{Acronyms}        & Acronyms \\
    Glossary & \texttt{Subscripts}      & Subscripts \\
    Lists    & \texttt{ListOfExamples}  & List of Examples \\
    Lists    & \texttt{FurtherReadingTitle} & Further Reading \\
    Science  & \texttt{Constant}        & const \\
    Science  & \texttt{PhysicalQuantity} & Physical Quantity \\
    Science  & \texttt{AdaptedFrom}     & Adapted from \\
    Reaction & \texttt{Reaction}        & Reaction \\
    Table    & \texttt{Contfoot}        & Continued on next page \\
    Table    & \texttt{Conthead}        & \ (Continued) \\
    Layout   & \texttt{BlankPage}       & Rest of this page intentionally left blank. \\
    \bottomrule
  \end{tabular}
\end{table}
